\section{Introduction}\label{sec:intro}

Volatility is persistent. Shocks to the variance of stock returns decay
slowly, the autocorrelations of squared and absolute returns stay positive
for months, and the log-periodogram of realized variance rises steeply
toward the origin \citep{ding1993long, andersen2003modeling}. A large
literature models this with fractional integration of order $d$, and an
equally large applied literature estimates $d$, or the related Hurst
exponent, over rolling windows and reads the resulting series as a
time-varying property of the market: persistence that strengthens in
crises, efficiency that improves or deteriorates, or adaptive behaviour in
the sense of \citet{lo2004adaptive}. Examples range from early work on
the Hurst exponent over time \citep{cajueiro2004hurst} to recent studies of
volatility persistence after COVID-19 \citep{veravaldes2022persistence} and
of the pandemic's effect on market efficiency \citep{takaishi2025}.

If rolling persistence is a genuine state of the market, it should carry
information about future volatility. This paper tests that implication and,
when it fails, asks why. We pose two questions. First, does a persistence
state built from rolling long-memory estimates forecast stock volatility
beyond what implied volatility already contains? Second, what does the
rolling persistence state measure?

We answer both on a panel of 115 S\&P~500 stocks from November 2001 to April
2026. Volatility is the daily Parkinson range variance, and the forecast
target is the log of its mean over the next 1, 5 or 22 trading days. The
benchmarks are the heterogeneous autoregressive (HAR) model of
\citet{corsi2009simple} and a HAR model augmented with the closing levels of
the VIX and MOVE indices (HAR-X). The persistence state enters through
own-stock rolling estimates of $d$, sector and cross-sectional averages of
those estimates, and their interactions with implied volatility. Every
forecast is point-in-time: predictors use information available at the close
of the forecast origin, and a training row enters a regression only if its
target window has closed by the origin. Each specification and its decision
rule were recorded in a dated experiment ledger before its results were
computed. Inference combines the panel Diebold--Mariano test with the
\citet{harvey1997testing} correction, the \citet{clark2007approximately} test
for nested models and the conditional predictive ability test of
\citet{giacomini2006tests}.

The answer to the first question is no. Implied volatility improves on HAR
by 3.1\%, 5.5\% and 1.8\% in mean squared error at the three horizons.
Persistence adds nothing on top. The best gain over HAR-X is 0.26\% at one
week, with a Diebold--Mariano statistic of 0.41, and the full model is
significantly worse than HAR-X at one day and one month. Pooled estimation
with stock fixed effects, a target that measures the slope of the variance
term structure, a single market-level series, and five machine-learning
estimators all give the same answer. The Clark--West test rejects the null
that the persistence terms have zero population coefficients, but the
Giacomini--White test never finds that the estimated models forecast better:
the information exists in population and cannot be extracted in practice.

The answer to the second question explains the first. The cross-sectional
mean of rolling $d$ estimates behaves as an indicator of whether a large,
long volatility episode lies inside the estimation window. It jumps when
such an episode enters the window and falls when the episode leaves, and the
exit date moves one-for-one with the window length: for windows of 500, 750
and 1{,}000 trading days the 2020 plateau ends on 11 March 2022, 10 March
2023 and 8 March 2024, each time with a fall in the bottom 0.2\% of all
weekly changes. A short-memory ARMA model fitted to calm-period data, with the
observed 2020 episode added, reproduces the calm level of the estimate, the
plateau and the exit. The largest VIX inside the window explains 81--83\% of
the variation in the state, against 18--28\% for the current VIX. A state
that is a lagged summary of past extreme episodes has little to add to an
implied-volatility index that already prices them.

A third result is practical. If HAR inputs end one day before the forecast
origin while the VIX is the origin-day close, HAR is handicapped and the
VIX's gain is overstated by about a third. With aligned inputs HAR's error
falls by 8.1\% at the daily horizon and the VIX's gain over HAR falls from
5.1\% to 3.1\%.

These results connect to three literatures. The first shows that structural
breaks, level shifts and regime switching produce spurious evidence of long
memory \citep{diebold2001long, granger2004occasional,
mikosch2004nonstationarities, smith2005level, perron2010long}, with tests and
estimators designed to tell the two apart \citep{qu2011test,
hou2014modified} and evidence that little genuine long memory survives once
level shifts are modelled \citep{varneskov2018combining}. That literature
concerns full-sample estimates. We document its time-series consequence for
rolling estimates, which are what applied work interprets, and give a simple
check: vary the window and see whether the turning points move with it. The
problem is old in other sciences. Mean shifts were shown to mimic the Hurst
phenomenon in river flows \citep{klemes1974hurst}, the separation of extreme
events from persistence goes back to the ``Noah'' and ``Joseph'' effects of
\citet{mandelbrot1968noah}, patchy composition mimics long-range correlation
in DNA \citep{karlin1993patchiness}, and rolling-window indicators of
critical slowing down and of dynamic brain connectivity have faced the same
critique \citep{dakos2008slowing, ditlevsen2010tipping,
boettiger2012quantifying, leonardi2015spurious}.

The second literature compares implied and time-series forecasts of
volatility \citep{christensen1998relation, fleming1998quality,
blair2001forecasting, busch2011role} and documents a variance risk premium
that varies with conditions \citep{bates2000post, chernov2007role,
bekaert2014vix, andersen2015risk}. Our timing result adds a concrete pitfall
to that comparison. The third literature asks whether volatility is rough
\citep{gatheral2018volatility, fukasawa2022consistent, cont2024rough} and
whether volatility forecasts create value in managed portfolios
\citep{moreira2017volatility, cederburg2020volatility,
barroso2021volatility}. We report both briefly: a roughness estimate of 0.06
on daily range data is reproduced by a non-rough model with measurement
noise, and managed portfolios built on any of our forecasts have the same
Sharpe ratio as the equal-weighted portfolio once returns are measured in
excess of the bill rate.

The paper proceeds as follows. Section~\ref{sec:data} describes the data
and the full-sample memory estimates. Section~\ref{sec:window} shows what
the rolling persistence state measures. Section~\ref{sec:design} sets out
the forecasting design and Section~\ref{sec:results} the results.
Section~\ref{sec:robust} reports robustness checks and the portfolio
exercise, and Section~\ref{sec:discussion} discusses implications and
limitations. The appendices collect supplementary results and the
registration record.
